\documentclass[../main.tex]{subfiles}
\ifSubfilesClassLoaded{\setcounter{chapter}{2}}{}
\begin{document}

\chapter{Pandas: Membaca CSV/Excel; \texttt{head}, \texttt{info}, \texttt{describe}; Tipe Data}

\begin{subcpmk}
  \item \textbf{Sub-CPMK-P2:} Memuat data tabular dan menginspeksi struktur, tipe, serta kualitas awal dengan Pandas (RPS Mg.\ 2).
\end{subcpmk}

\SesiRPS{2}{P2}{$\sim$170 menit}{CPMK-P2}

\section{Kemampuan akhir sesi}

Memuat CSV dan Excel; menggunakan \texttt{head}, \texttt{tail}, \texttt{info}, \texttt{describe}, \texttt{dtypes}; mengajukan hipotesis awal kolom target vs fitur.

\section{Teori dan bahan kajian}

\begin{konsep}
\texttt{DataFrame} adalah tabel dengan indeks baris dan nama kolom. \texttt{dtypes} membedakan \texttt{int64}, \texttt{float64}, \texttt{object} (sering berarti teks atau campuran). Fungsi \texttt{info()} memberi ringkasan non-null dan tipe; \texttt{describe()} merangkum statistik numerik secara default.
\end{konsep}

\begin{catatan}
Path berkas: gunakan path relatif ke folder kerja notebook atau unggah ke Colab dengan \texttt{files.upload()}. Dokumentasikan path di Markdown agar asisten dapat mereproduksi.
\end{catatan}

\section{K3 dan etika}

Dataset hanya dari sumber yang diizinkan; jangan menyematkan token API atau path absolut sensitif di notebook yang diunggah ke LMS.

\section{Sarana}

\texttt{pandas}, \texttt{numpy} (dependensi umum), \texttt{openpyxl} untuk \texttt{.xlsx}.

\section{Prosedur kerja}

\begin{enumerate}[leftmargin=*]
  \item \texttt{import pandas as pd}
  \item Muat satu CSV (\texttt{read\_csv}); tampilkan \texttt{head(10)} dan \texttt{info()}.
  \item Muat sumber kedua (CSV lain atau Excel); bandingkan \texttt{shape} dan \texttt{dtypes}.
  \item Di Markdown, tulis kolom kandidat \textbf{target} (regresi/klasifikasi) dan alasan singkat.
  \item Sinkronkan dengan \texttt{notebooks/minggu\_02.ipynb}.
\end{enumerate}

\section{Contoh kode}

\begin{lstlisting}[caption=Import dan muat CSV]
import pandas as pd

# Ganti path dengan file Anda (contoh: dataset publik)
df = pd.read_csv("data/contoh.csv")  # encoding="utf-8" bila perlu
print(df.shape)
print(df.head())
\end{lstlisting}

\begin{lstlisting}[caption=Ringkasan struktur dan statistik]
print(df.info())
print(df.describe(include="all").T.head(12))
\end{lstlisting}

\begin{lstlisting}[caption=Membaca Excel (satu sheet)]
df_xlsx = pd.read_excel("data/contoh.xlsx", sheet_name=0)
print(df_xlsx.dtypes.value_counts())
\end{lstlisting}

\begin{lstlisting}[caption=Membandingkan dua kerangka data]
def ringkas(nama, frame):
    return {
        "nama": nama,
        "baris": len(frame),
        "kolom": frame.shape[1],
        "missing_total": frame.isna().sum().sum(),
    }

print(ringkas("csv", df))
print(ringkas("xlsx", df_xlsx))
\end{lstlisting}

\section{Referensi sesi}

\begin{itemize}[leftmargin=*]
  \item pandas User Guide IO. \url{https://pandas.pydata.org/docs/user_guide/io.html}
  \item Kaggle Learn Pandas. \url{https://www.kaggle.com/learn/pandas}
  \item Microsoft Learn Pandas. \url{https://learn.microsoft.com/en-us/training/modules/pandas-data-science/}
\end{itemize}

\begin{aktivitas}
  \item Muat dua dataset berbeda; dokumentasikan URL dan lisensi.
  \item Buat tabel Markdown ringkas: nama kolom, tipe, persen missing (hitung manual atau dengan kode).
\end{aktivitas}

\begin{latihan}
  \item Kapan Anda memilih \texttt{read\_csv(..., sep=";")}?
  \item Apa implikasi kolom numerik yang tersimpan sebagai \texttt{object}?
\end{latihan}

\begin{asesmen}
\textbf{Tugas modul:} \texttt{read\_*} tanpa error; dokumentasi path/URL; hipotesis target/fitur.

\textbf{Siap lanjut:} memahami \texttt{info()} dan \texttt{describe()} untuk minggu cleaning.
\end{asesmen}

\begin{checklist}
  \item Saya dapat memuat CSV dan (opsional) Excel
  \item Saya membaca output \texttt{info} dan \texttt{describe}
\end{checklist}

\begin{rangkuman}
Minggu 2 menegaskan I/O Pandas dan inspeksi awal sebagai fondasi prapemrosesan.

\textbf{Kata kunci:} \texttt{read\_csv}, \texttt{read\_excel}, \texttt{DataFrame}, dtypes
\end{rangkuman}

\end{document}
